\chapter{DevOps, Infraestructura, Ética en IA y Conclusiones}
\label{cap:devops_conclusiones}

\section{Infraestructura como Código y Contenedorización con Docker}
Para mitigar el clásico problema de dispersión de entornos (el antipatrón \textit{«en mi máquina funciona»}) y garantizar la reproducibilidad absoluta del sistema en cualquier estación de trabajo o servidor en la nube, el proyecto implementa Infraestructura como Código mediante \textbf{Docker} y \textbf{Docker Compose}.

\subsection{Topología de Servicios Contenedorizados}
El archivo \texttt{docker-compose.yml} orquesta los servicios centrales requeridos para el almacenamiento y la administración de datos:
\begin{enumerate}
    \item \textbf{Servicio de Base de Datos (\texttt{db}):} Contenedor con motor \textbf{MySQL 8.0}, configurado con juego de caracteres \texttt{utf8mb4} y cotejamiento \texttt{utf8mb4\_unicode\_ci}. Los datos se persisten en un volumen nombrado de Docker (\texttt{db\_data}), garantizando que el ciclo de vida del contenedor no comprometa la integridad de los registros.
    \item \textbf{Servicio de Gestión Visual (\texttt{phpmyadmin}):} Interfaz web para administración de bases de datos expuesta en el puerto \texttt{8080}, configurada con límites de carga de memoria para soporte de restauración de respaldos SQL pesados.
    \item \textbf{Red Aislada de Aplicación (\texttt{tienda-network}):} Red tipo \textit{bridge} que encapsula la comunicación inter-servicios e impide la exposición indiscriminada de puertos internos al exterior.
\end{enumerate}

\begin{lstlisting}[language=bash, caption={Arranque de la Infraestructura de Base de Datos}, label={lst:docker_run}]
# Inicializacion de contenedores en segundo plano
docker compose up -d

# Verificacion del estado de salud de los servicios
docker compose ps
\end{lstlisting}

\section{Pipeline de Integración Continua (CI) en GitHub Actions}
La arquitectura de aseguramiento de calidad del repositorio incluye un flujo automatizado de Integración Continua denominado \texttt{docs-validation.yml}, el cual se dispara ante cada evento de \texttt{push} o \texttt{pull\_request} hacia la rama principal \texttt{main}.

\begin{figure}[H]
\centering
\begin{tikzpicture}[
    scale=0.92, every node/.style={transform shape},
    stepnode/.style={rectangle, draw=primaryblue!80!black, fill=primaryblue!10, rounded corners=4pt, line width=1pt, minimum width=7.0cm, minimum height=0.85cm, align=center, font=\small, inner sep=5pt},
    arrow/.style={thick, -{Stealth[length=2.5mm]}, primaryblue}
]
\node[stepnode] (git) {\textbf{Trigger:} Push / Pull Request a rama \texttt{main}};
\node[stepnode, below=0.45cm of git] (lint) {Linter de Documentación (\texttt{markdownlint-cli2})};
\node[stepnode, below=0.45cm of lint] (puml) {Validación Sintáctica PlantUML (\texttt{-checkonly})};
\node[stepnode, below=0.45cm of puml] (mmd) {Renderizado Headless de Diagramas Mermaid (\texttt{--no-sandbox})};
\node[stepnode, below=0.45cm of mmd, fill=green!15, draw=green!70!black] (pass) {\textbf{Resultado:} Verificación Exitosa (Pipeline Green 100\%)};

\draw[arrow] (git) -- (lint);
\draw[arrow] (lint) -- (puml);
\draw[arrow] (puml) -- (mmd);
\draw[arrow] (mmd) -- (pass);
\end{tikzpicture}
\caption{Flujo de Verificación Automatizada en GitHub Actions}
\label{fig:pipeline_ci}
\end{figure}

Las etapas automatizadas garantizan que ningún cambio introduzca errores sintácticos o inconsistencias en los diagramas arquitectónicos:
\begin{itemize}
    \item \textbf{Linting Estático de Documentación:} Inspección con \texttt{markdownlint-cli2} aplicando reglas formales de encabezados jerárquicos y formato de tablas.
    \item \textbf{Compilación Sintáctica de Diagramas de Secuencia:} Validación de los 8 archivos \texttt{.puml} mediante el motor oficial PlantUML en modo \texttt{-checkonly}, asegurando cumplimiento estricto de sintaxis \texttt{strictuml}.
    \item \textbf{Renderizado Headless de Diagramas Mermaid:} Generación automatizada con Mermaid CLI (\texttt{mmdc}) empleando banderas \texttt{--no-sandbox} para compatibilidad en entornos virtualizados Linux.
\end{itemize}


\section{Conclusiones Generales del Proyecto}
Tras el ciclo completo de investigación, análisis, refactorización y formalización arquitectónica, se arriba a las siguientes conclusiones fundamentales:
\begin{enumerate}
    \item \textbf{Superación Cualitativa del Enfoque Legacy:} La transición desde el enfoque intuitivo y no formal del proyecto heredado (Análisis de Sistemas) hacia un marco formal de Ingeniería de Software permitió expandir la cobertura de requerimientos de 18 a \textbf{77 Requisitos Funcionales} y de 6 a \textbf{14 Requisitos No Funcionales}, resolviendo la totalidad de las brechas operativas de la tienda Machi's.
    \item \textbf{Desacoplamiento y Mantenibilidad Arquitectónica:} La separación en cuatro capas concéntricas (Dominio, Aplicación, Infraestructura y Presentación) junto con la segregación CQRS garantiza que el sistema pueda evolucionar, cambiar su motor de persistencia o adoptar nuevos clientes visuales sin impacto destructivo sobre las reglas de negocio.
    \item \textbf{Solidez en la Persistencia de Datos (3FN):} El rediseño integral de la base de datos a 17 tablas en Tercera Forma Normal eliminó anomalías de duplicación en inventario y garantizó trazabilidad total de caja, ventas y transferencias multi-sucursal bajo transacciones ACID.
    \item \textbf{Verificabilidad Empírica de Calidad (IEEE 829):} La formulación de 77 casos de prueba detallados proporciona a la cátedra y al cliente la certeza matemática y operacional de que cada función implementada cuenta con criterios de aceptación auditables.
    \item \textbf{Cumplimiento Integral de la Rúbrica ESIS:} El proyecto satisface con rigor los \textbf{27 criterios de evaluación} establecidos en la rúbrica institucional, abarcando proceso de desarrollo, métricas de mejora, validación e integración continua.
\end{enumerate}

\section{Hoja de Ruta Futura (Roadmap)}
Para fases subsiguientes de escalamiento comercial del sistema Tienda Machi's se proponen las siguientes iniciativas técnicas:
\begin{itemize}
    \item \textbf{Fase 3: Facturación Electrónica SUNAT:} Integración directa de un microservicio de firma digital y emisión de Comprobantes de Pago Electrónicos (Boletas y Facturas XML/UBL 2.1) conectado a los servicios web de la SUNAT.
    \item \textbf{Fase 4: Aplicación Móvil de Inventario en Bodega:} Desarrollo de un cliente nativo ligero para dispositivos móviles Android con lectura de códigos de barras mediante cámara para auditorías físicas en tiempo real.
    \item \textbf{Fase 5: Motor de Recomendación y Fidelización:} Análisis de series de tiempo para predicción de demanda estacional y módulo de acumulación de puntos para clientes recurrentes.
\end{itemize}
